\section{Vertex gauges and holonomy}%
\label{sec:peps_regular_gauges_holonomy}

Group elements inserted on the internal bonds of a region can be moved by
gauge transformations at its vertices. Along a spanning tree they can be
removed, and what remains is the holonomy of the inserted elements along
closed walks.

\begin{lemma}[Twisted canonical block in spanning-tree coordinates]%
    \label{thm:peps_regular_twisted_projector_coordinates}
    \lean{TNLean.PEPS.regularProjectorTwistedRegionMatrix_coordinates}
    \leanok
    \uses{def:peps_regular_region_coordinates,
        thm:peps_regular_region_tree_gauge,
        def:peps_regular_projector_twisted_region}
    Choose a spanning tree $T$ rooted at $o$ in the finite region $R$,
    and insert arbitrary operators $u_e\in G$ on its incident bonds.
    Let $k_o=1$, $k_tk_s^{-1}=u_e$ on $T$, and
    $\omega_e=k_t^{-1}u_ek_s$ on the internal edges.
    Let $\theta'$ be the crossing configuration transformed by
    (\ref{eq:peps_tree_boundary_transport}). In the original physical
    coordinates $\alpha(y,a,r,z)$, the twisted canonical matrix satisfies
    \begin{align}
        M_R^u\bigl(\alpha(y,a,r,z),\theta\bigr)
         & = |G|^{-|R|}\sum_{x\in G}
        \mathbf 1_{\{y=x\theta'\}}
        \mathbf 1_{\{z_e=x\omega_ex^{-1}\ (e\in I\setminus T)\}}.
        \label{eq:peps_twisted_projector_coordinates}
    \end{align}
    The reference labels $a$ and rooted vertex labels $r$ do not affect
    the coefficient. Operators on crossing bonds are allowed, and no
    condition is imposed on the internal residuals or the size of the
    boundary. This is a finite-contraction form of the regular
    coordinate changes in \cite[\S6.2--6.3]{Schuch2010PEPS}.
\end{lemma}
\begin{proof}\leanok
    Expanding the original contraction into its internal bond labels
    $\eta$ and local translations $q$, and setting
    $q'_v=q_vk_v$ and $\eta'_e=k_s^{-1}\eta_e$ on the tree,
    rewrites the endpoint equations as
    $r_sa_e=q'_s\eta'_e$ and $r_ta_e=q'_t\eta'_e$.
    Rooted tree-gradient uniqueness yields $q'_v=r_vx$, with $x=q_o$.
    Thus $q_v=r_vxk_v^{-1}$, and internal references are uniquely fixed
    as $\eta_e=k_sx^{-1}a_e$. On each remaining internal edge the
    head equation is equivalent to $z_e=x\omega_ex^{-1}$.

    At a crossing edge, the unique reference equals
    $k_sx^{-1}y_e$ when $s\in R$, and
    $u_e^{-1}k_tx^{-1}y_e$ when $t\in R$.
    Fixing these references to $\theta$ is exactly $y=x\theta'$.
    Conversely, these boundary and cycle equations reconstruct a
    compatible pair $(\eta,q)$. Hence the compatible pairs are in
    bijection with the surviving $x\in G$, and the original local
    averages contribute the factor $|G|^{-|R|}$ in
    (\ref{eq:peps_twisted_projector_coordinates}).
\end{proof}

\begin{lemma}[Mixed Gram matrices of twisted canonical blocks]%
    \label{thm:peps_twisted_projector_cross_gram_coordinates}
    \lean{TNLean.PEPS.regularProjectorTwistedRegionMatrix_crossGram_coordinates,
        TNLean.PEPS.regularProjectorTwistedRegionMatrix_crossGram_eq_zero_of_cycles_not_conjugate}
    \leanok
    Fix a region $R$, a spanning tree, and a root, and let $u,w$ be
    arbitrary bond insertions, including crossing operators. Write
    $\omega^u,\omega^w$ for their cycle residual tuples and
    $b_u,b_w$ for their respective boundary maps. Let $N$ be
    the number of pairs $(a,r)$ of internal references and vertex
    labels normalized at the root. The canonical matrices obey
    \begin{align}
        \bigl((M_R^u)^\dagger M_R^w\bigr)(\theta,\vartheta)
         & = |G|^{-2|R|}N\sum_{x,t\in G}
        \mathbf 1_{\{x b_u(\theta)=t b_w(\vartheta)\}}
        \mathbf 1_{\{x\omega^ux^{-1}=t\omega^wt^{-1}\}}.
        \label{eq:peps_twisted_projector_cross_gram}
    \end{align}
    In particular, if no $s\in G$ satisfies
    $\omega^u=s\omega^ws^{-1}$ coordinatewise, then
    $(M_R^u)^\dagger M_R^w=0$. Thus these two blocks have
    orthogonal physical ranges. The boundary may be empty; neither
    flatness nor a mixed Gram identity is assumed.
    This is a finite-region comparison supporting the coherent closure
    arguments in \cite[\S6.3--6.4]{Schuch2010PEPS}.
\end{lemma}
\begin{proof}\leanok
    \uses{thm:peps_regular_twisted_projector_coordinates}
    Reindexing the physical row sum by the reversible spanning-tree
    coordinates and substituting
    (\ref{eq:peps_twisted_projector_coordinates}) for each coefficient
    leaves the factors indexed by $(a,r)$ constant, contributing $N$.
    Summing over the shared boundary and cycle coordinates pairs the
    two indicators, leaving precisely the equalities in
    (\ref{eq:peps_twisted_projector_cross_gram}).
    Any surviving cycle equality implies
    $\omega^u=(x^{-1}t)\omega^w(x^{-1}t)^{-1}$.
    If simultaneous conjugacy fails, every summand vanishes.
\end{proof}

\begin{lemma}[Vertex gauges in a regular-region contraction]%
    \label{thm:peps_regular_region_gauge_contraction}
    \lean{TNLean.PEPS.regularRegionGaugeIncidentLabels_boundary,
        TNLean.PEPS.regularRegionGaugeEdgeOperators,
        TNLean.PEPS.regularRegionGaugeIncidentLabels,
        TNLean.PEPS.regularRegionGaugeIncidentLabels_twisted,
        TNLean.PEPS.openRegionWeight_regularRegionGauge,
        TNLean.PEPS.openRegionWeight_eq_untwisted_of_regularRegionGaugeResidual_eq_one}
    \leanok
    \uses{def:peps_regular_twisted_region, def:peps_open_region_contraction}
    Suppose each original site coefficient is invariant under simultaneous
    left translation of its incident regular labels. Let $k_v\in G$ be
    arbitrary gauges at the vertices of a finite region $R$, and let $u_e$
    be arbitrary inserted operators, including those on crossing bonds.
    Put $\widetilde u_e=k_t^{-1}u_ek_s$ on internal edges $e=(s,t)$,
    and put $\widetilde u_e=1$ on all other edges. With the transformed
    crossing configuration $\theta'$ of
    (\ref{eq:peps_tree_boundary_transport}), the physical columns satisfy
    \begin{align}
        T_R^u(\theta) & =T_R^{\widetilde u}(\theta').
    \end{align}
    If every internal residual is the identity, this reduces to
    $T_R^u(\theta)=T_R(\theta')$ for the original untwisted region.
    No connectedness or Gram identity is required. These identities
    implement the regular-coordinate changes in
    \cite[\S6.2--6.3]{Schuch2010PEPS} inside the original finite contraction.
\end{lemma}

\begin{proof}\leanok
    \uses{thm:peps_regular_region_tree_gauge}
    Reindexing the sum over incident edge references multiplies the reference
    by $k_s^{-1}$ on an internal edge and applies
    the map in (\ref{eq:peps_tree_boundary_transport}) on a crossing edge. Each change
    is a bijection. Its boundary restriction is exactly $\theta\mapsto\theta'$.
    At every region vertex, the resulting twisted half-edge labels are
    $k_v^{-1}$ times the original twisted labels. Local simultaneous
    invariance therefore leaves each site coefficient unchanged, and
    hence leaves their product and the contracted sum unchanged.
\end{proof}

\begin{lemma}[The transformed boundary projector]%
    \label{thm:peps_regular_region_gauge_boundary}
    \lean{TNLean.PEPS.regularRegionBoundaryGaugeEquiv_apply,
        TNLean.PEPS.regularRegionBoundaryGaugeEquiv,
        TNLean.PEPS.regularRegionGaugeBoundaryProjector,
        TNLean.PEPS.regularRegionGaugeBoundaryProjector_mul_self,
        TNLean.PEPS.regularRegionGaugeBoundaryProjector_conjTranspose,
        TNLean.PEPS.rank_regularRegionGaugeBoundaryProjector,
        TNLean.PEPS.trace_regularRegionGaugeBoundaryProjector}
    \leanok
    \uses{thm:peps_regular_region_tree_gauge}
    Enumerate the crossing bonds of a region by $\{1,\ldots,b\}$, and let
    $F$ be the resulting permutation of boundary configurations induced
    by the vertex gauge and the crossing operators. The pulled-back
    averaging projector is
    \begin{align}
        \Pi_{\partial R}^{k,u}(x,y)=\Pi_{\partial R}(F(x),F(y)).
    \end{align}
    It is idempotent. If $b\geq1$, it is self-adjoint, has rank
    $|G|^{b-1}$, and has trace $|G|^{b-1}$. This coordinate change is
    part of the boundary disentangling argument of
    \cite[Theorem~6.9]{Schuch2010PEPS}; it imposes no condition on
    internal twists.
\end{lemma}

\begin{proof}\leanok
    Reindexing the common summation variable in the square of the matrix
    by $F$ derives idempotence from that of the original averaging
    projector. Taking the adjoint commutes with applying the same
    permutation to both indices. Finally, a basis permutation preserves
    rank, and reindexing the diagonal sum preserves trace.
\end{proof}

\begin{lemma}[The Gram matrix with removable internal operators]%
    \label{thm:peps_regular_region_gauge_gram}
    \lean{TNLean.PEPS.regularTwistedOpenRegionMatrix_regularRegionGauge,
        TNLean.PEPS.regularTwistedOpenRegionMatrix_eq_submatrix_of_regularRegionGaugeResidual_eq_one,
        TNLean.PEPS.exists_positive_gram_regularTwistedOpenRegionMatrix_of_regularRegionGaugeResidual_eq_one,
        TNLean.PEPS.rank_regularTwistedOpenRegionMatrix_of_regularRegionGaugeResidual_eq_one}
    \leanok
    Number the crossing bonds and write $F$ for the boundary permutation
    induced by an arbitrary vertex gauge $k$ and inserted operators $u$.
    If the original sites have simultaneous regular invariance, the
    ordered coefficient matrices satisfy
    \begin{align}
        M_R^u(\sigma,x) & =M_R^{\widetilde u}(\sigma,F(x)).
    \end{align}
    Suppose further that $R$ is connected, the sites are $G$-isometric,
    and every internal residual $k_t^{-1}u_ek_s$ is the identity.
    Then $M_R^u(\sigma,x)=M_R(\sigma,F(x))$ and there is a positive
    scalar $c_R$ such that
    \begin{align}
        (M_R^u)^\dagger M_R^u & =c_R\Pi_{\partial R}^{k,u}.
    \end{align}
    If the boundary contains $b>0$ bonds, the matrix has rank
    $|G|^{b-1}$. The Gram statement allows an empty boundary and derives the block Gram
    identity from the local sites. It is the removable-operator case of
    the boundary-isometry argument in \cite[Theorem~6.9]{Schuch2010PEPS},
    not the unrestricted disk theorem
    \cite{gap:rmp_peps_quantum_double_g_isometry}.
\end{lemma}

\begin{proof}\leanok
    \uses{thm:peps_regular_region_gauge_contraction, thm:peps_regular_region_gauge_boundary, thm:peps_connected_regular_region_isometry}
    The ordered coefficient identities follow from the contraction
    identity in Lemma~\ref{thm:peps_regular_region_gauge_contraction}.
    With identity internal residuals, the matrix is the untwisted block
    with its columns permuted by $F$. Local regular isometry and
    connectedness give $M_R^\dagger M_R=c_R\Pi_{\partial R}$ with $c_R>0$.
    Pulling both boundary indices back along $F$ therefore gives the
    claimed Gram matrix. Its rank equals the rank of the matrix
    because $\operatorname{rank}(M^\dagger M)=\operatorname{rank}(M)$;
    the positive scalar does not change the rank of the boundary projector.
\end{proof}

\begin{lemma}[Vertex gauges of a closed regular PEPS]%
    \label{thm:peps_regular_closed_gauge}
    \lean{TNLean.PEPS.sameState_regularVertexGauge,
        TNLean.PEPS.regularVertexGaugeOperators}
    \leanok
    \uses{thm:peps_regular_region_gauge_contraction}
    Suppose the site coefficients $a_v$ are invariant under simultaneous left translation of their incident group labels. For vertex gauges $k_v\in G$, replacing the bond operators $u_e$ on $e=(s,t)$ by
    \begin{align}
        u_e' & = k_t^{-1}u_e k_s
        \label{eq:peps_regular_closed_gauge}
    \end{align}
    leaves every closed physical PEPS coefficient unchanged (without requiring connectedness or isometry \cite[\S5.3]{Schuch2010PEPS}).
\end{lemma}

\begin{proof}\leanok
    Reindexing each bond label $\eta_e\mapsto k_s^{-1}\eta_e$ on $e=(s,t)$ sends the tail half-edge label to $k_s^{-1}\eta_e$ and the head half-edge label to $u_e'k_s^{-1}\eta_e=k_t^{-1}u_e\eta_e$, left-translating all incident labels at each vertex $v$ by $k_v^{-1}$.
\end{proof}

\begin{lemma}[Nonabelian holonomy along native walks]%
    \label{thm:peps_regular_walk_holonomy}
    \lean{TNLean.PEPS.regularDirectedTransport_of_lt,
        TNLean.PEPS.regularDirectedTransport,
        TNLean.PEPS.regularDirectedTransport_of_gt,
        TNLean.PEPS.regularDirectedTransport_symm,
        TNLean.PEPS.regularDirectedTransport_gauge,
        TNLean.PEPS.regularWalkHolonomy,
        TNLean.PEPS.regularWalkHolonomy_nil,
        TNLean.PEPS.regularWalkHolonomy_cons,
        TNLean.PEPS.regularWalkHolonomy_append,
        TNLean.PEPS.regularWalkHolonomy_reverse,
        TNLean.PEPS.regularWalkHolonomy_gauge,
        TNLean.PEPS.regularWalkHolonomy_gauge_loop,
        TNLean.PEPS.regularWalkHolonomy_gradient,
        TNLean.PEPS.regularWalkHolonomy_gradient_loop}
    \leanok
    \uses{thm:peps_regular_closed_gauge}
    Let $G$ be any group, and assign $u_e\in G$ to each edge of a graph
    with ordered vertices. A directed step along $e=(s,t)$ uses $u_e$
    from $s$ to $t$ and $u_e^{-1}$ in the reverse direction. For a walk
    $p$, let $H(u,p)$ be the product of these operators in reverse
    traversal order; the empty walk has product one. If $p:v\to w$
    and $q:w\to z$, and if $u'_e=k_t^{-1}u_ek_s$, then
    \begin{align}
        H(u,pq)     & =H(u,q)H(u,p),      \\
        H(u,p^{-1}) & =H(u,p)^{-1},       \\
        H(u',p)     & =k_w^{-1}H(u,p)k_v.
        \label{eq:peps_walk_gauge}
    \end{align}
    In particular a closed-walk holonomy transforms by conjugation at
    its base vertex. A gradient assignment $u_e=k_tk_s^{-1}$ has
    $H(u,p)=k_wk_v^{-1}$ and hence identity holonomy on closed walks.
    These are group identities underlying the seam deformation and
    accessible complement cycle calculations in
    \cite[\S5.3 and Theorem~6.9]{Schuch2010PEPS}; no boundary-word
    relation or tensor factorization is assumed.
\end{lemma}

\begin{proof}\leanok
    Reversing a directed edge takes its inverse. The concatenation and
    reversal formulas follow by induction on the walk. The directed
    gauged operator from $v$ to $w$ is $k_w^{-1}u_{v\to w}k_v$.
    In the reverse-ordered product, adjacent endpoint gauges cancel,
    leaving (\ref{eq:peps_walk_gauge}). For a gradient, the factors
    telescope directly to $k_wk_v^{-1}$.
\end{proof}

\begin{lemma}[Closure holonomy from oriented seam crossing numbers]%
    \label{thm:peps_torus_walk_winding}
    \lean{TNLean.PEPS.torusDirectedWinding_right,
        TNLean.PEPS.torusEdgeWinding,
        TNLean.PEPS.torusDirectedWinding,
        TNLean.PEPS.torusWalkWinding,
        TNLean.PEPS.torusDirectedWinding_up,
        TNLean.PEPS.torusDirectedWinding_symm,
        TNLean.PEPS.torusWalkWinding_append,
        TNLean.PEPS.torusWalkWinding_reverse,
        TNLean.PEPS.regularWalkHolonomy_torusClosureEdgeAssignment,
        TNLean.PEPS.regularWalkHolonomy_torusClosure_eq_of_winding_eq}
    \leanok
    \uses{thm:peps_regular_walk_holonomy, thm:peps_regular_torus_closure_cut_identification}
    Consider the native torus graph with both periods at least three.
    For a directed step, count $+1$ at a rightward horizontal seam
    crossing and at an upward vertical seam crossing, count $-1$ in
    the respective reverse directions, and count zero elsewhere.
    Let $N(p)=(n(p),m(p))\in\mathbb Z^2$ be the sum of these pairs
    along a walk $p$. Concatenation and reversal satisfy
    $N(pq)=N(p)+N(q)$ and $N(p^{-1})=-N(p)$.

    Let $g,h\in G$ commute, and let $u(g,h)$ be the ordered-edge
    closure assignment. Then
    \begin{align}
        H(u(g,h),p)=h^{n(p)}(g^{-1})^{m(p)}.
        \label{eq:peps_winding_holonomy}
    \end{align}
    Consequently any two walks with equal seam crossing numbers have
    equal closure holonomies, even when their endpoints differ.
    This is the walk identity used in the seam deformation and
    complement construction of
    \cite[\S5.3 and proof of Theorem~6.9]{Schuch2010PEPS}.
\end{lemma}

\begin{proof}\leanok
    Both native seam orientations are opposite to their ordered-edge
    orientations. Thus the rightward seam operator is $h$ and the
    upward seam operator is $g^{-1}$. Reversing a directed step
    negates its crossing numbers and inverts its operator. Every step
    therefore satisfies~(\ref{eq:peps_winding_holonomy}). Induction on
    the walk adds the integer exponents; commutativity of $g$ and $h$
    permits the two powers to be grouped. The addition and reversal
    rules for $N$ follow by induction from the corresponding rules for
    directed steps.
\end{proof}

\begin{lemma}[Walk words in the region cycle labels]%
    \label{thm:peps_regular_region_cycle_holonomy}
    \lean{TNLean.PEPS.regularRegionCycleWord_append,
        TNLean.PEPS.RegionCycleEdge,
        TNLean.PEPS.regularRegionInternalOperators,
        TNLean.PEPS.regularRegionCycleOperators,
        TNLean.PEPS.regularRegionCycleWord,
        TNLean.PEPS.regularRegionCycleWord_reverse,
        TNLean.PEPS.regularRegionCycleWord_eval,
        TNLean.PEPS.regularRegionCycleWord_eval_conj,
        TNLean.PEPS.regularRegionCycleOperators_treeGauge,
        TNLean.PEPS.regularRegionCycleWord_eval_treeResidual,
        TNLean.PEPS.regularRegionCycleWord_eval_rootLoop,
        TNLean.PEPS.regularWalkHolonomy_eq_cycleWord_eval}
    \leanok
    Let $T$ be a chosen spanning tree of an induced region graph, rooted
    at $o$, and let $C$ be its native internal edges outside $T$.
    For a region walk $p$, construct a word $W_T(p)$ in the
    free group on $C$: traverse $p$, use the generator or its inverse
    on each non-tree edge according to its orientation, omit tree edges,
    and multiply in reverse traversal order. Given any tuple $z\in G^C$,
    write $\operatorname{ev}_z$ for evaluation of this word. Assigning
    identity operators to tree edges and $z_e$ to the remaining edges
    gives walk holonomy equal to $\operatorname{ev}_zW_T(p)$.
    Concatenation satisfies $W_T(pq)=W_T(q)W_T(p)$ and reversal satisfies
    $W_T(p^{-1})=W_T(p)^{-1}$.

    For arbitrary native bond operators $u$, let $k$ be the existing
    root-normalized tree gauge and $\omega$ its cycle residual tuple.
    Its reduced assignment is precisely identity on the tree and
    $\omega_e$ on $C$. For $p:v\to w$ one therefore has
    \begin{align}
        \operatorname{ev}_{\omega}W_T(p) & =k_w^{-1}H(u,p)k_v,                           \\
        H(u,p)                           & =k_w\operatorname{ev}_{\omega}W_T(p)k_v^{-1}, \\
        \operatorname{ev}_{xzx^{-1}}W_T(p)
                                         & =x\operatorname{ev}_zW_T(p)x^{-1}.
    \end{align}
    For a loop based at $o$, normalization gives
    $\operatorname{ev}_{\omega}W_T(p)=H(u,p)$.
    Word construction and conjugation require no finiteness or
    commutativity; the existing normalized tree gauge is used for finite $G$.
    These are walk-word identities for the accessible complement
    cycle argument in \cite[Theorem~6.9]{Schuch2010PEPS}. They do not
    assert existence of walks representing a prescribed boundary map.
\end{lemma}

\begin{proof}\leanok
    \uses{thm:peps_regular_walk_holonomy, thm:peps_regular_region_tree_gauge}
    Applying the native walk product to the free generators and using the fact that evaluation
    commutes with each product and inverse yields the residual holonomy.
    The tree gauge removes each tree operator and retains exactly the
    prescribed cycle residual on every other edge. Endpoint gauge
    covariance gives the first two identities. Conjugating every
    residual operator conjugates the whole ordered product. At the
    normalized root the endpoint gauges are both the identity.
\end{proof}
